% ================================================
% 文件: 07_监测接收机.tex
% 章节: 第12章 监测接收机
% 所属部分: 接收机技术
% 版本: v1.0
% ================================================
% 本章目录结构（树状）:
% \chapter{监测接收机}
% ├── \section{监测接收机概述}
% ├── \section{监测接收机基本原理}
% ├── \section{监测接收机技术指标}
% ├── \section{监测接收机设计要点}
% └── \section{监测接收机应用场景}
% ================================================

\chapter{监测接收机}
\label{chap:monitoring_receiver}

\section{监测接收机概述}
\label{sec:monitoring_receiver_intro}

监测接收机是用于监测无线电频谱使用情况的专业设备。
它能够扫描、
分析和记录无线电信号，
用于频谱管理、
干扰监测和信号分析。

监测接收机的主要特点：
\begin{itemize}
    \item \textbf{宽频段覆盖}：通常覆盖从低频到微波的宽范围
    \item \textbf{高灵敏度}：能够检测微弱信号
    \item \textbf{实时分析}：能够实时分析信号特征
    \item \textbf{记录功能}：能够记录信号数据供后续分析
\end{itemize}

监测接收机（亦称频谱监测接收机或监测测向接收机）服务于无线电管理、
电磁频谱管控与信号情报等领域，
其工作对象不是某个“预定用户”的信号，
而是整个频段内所有可能存在的、
已知或未知的信号。
因此它与通信接收机在目标上不同：
通信接收机追求“对某特定信号的可靠解调”，
监测接收机追求“对全频段信号的发现、
刻画与定位”。
这决定了它必须具备极宽的频段覆盖、
良好的幅度/频率测量精度、
足够高的动态范围，
以及把频谱以“全景”方式呈现并长期记录的能力。

现代监测接收机已高度数字化：
射频前端完成变频与放大后，
信号在数字中频（Digital IF）上以高速 A/D 采样，
所有扫描、
检波、
FFT 谱分析、
调制识别与测向解算均在处理器或 FPGA/DSP 中完成。
配合数据库与地理信息系统（GIS），
监测站可对台站是否合规、
频段是否被占用、
干扰来自何方给出定量结论。
后续各节依次说明其原理、
指标、
设计与应用。

监测接收机与通信接收机目标迥异：
通信接收机追求“对某个预定信号的可靠解调”，
监测接收机追求“对全频段所有信号的发现、
测量、
刻画与定位”，
对象可能是未知、
瞬现甚至非法的信号。
因此它不是以“解调质量”为首要判据，
而以“发现概率、
频率/幅度测量精度、
动态范围与定位精度”为核心能力。
单个监测站难以覆盖全域，
实际系统常以分布式监测网组成，
各站把频谱轨迹、
信号参数与地理位置上报中心，
统一做态势呈现与干扰定位。

频谱占用度统计是监测的基本产出：
在一段时间窗内对每频段统计“被占用的时间占比”，
据此评估频率规划是否合理、
某频带是否闲置可重新指配。
结合 GIS 与台站数据库，
可将监测结果叠加到地图上，
直观显示“谁在何处发射、
是否符合执照”。
法规层面，
许多国家设有无线电管理机构的监测职责，
监测数据是频谱指配、
 
interference 查处与设备型号核准的技术依据。
理解这一“网络 + 数据库 + 地理信息”的运行框架，
有助于理解监测接收机为何强调宽覆盖、
可记录与可远程控制。

监测站按形态分为固定站、
移动监测车、
便携/手持设备与空间监测平台。
固定站常年值守重点频段与边境，
天线高、
灵敏度好；
监测车机动逼近干扰源；
便携设备用于现场查证；
空间平台则从轨道俯瞰大范围。
与通用频谱分析仪不同，
监测接收机强调“可程控、
可长时间连续记录、
可多站协同与可远程调度”，
并常内置数据库接口与 GIS，
输出的是带时间—位置—参数的结构化监测数据而非单幅截图。

数据记录格式影响后续分析能力：
除频谱轨迹外，
高级设备可记录 IQ 样本（原始复数基带），
便于离线重做解调、
调制识别与测向；
记录常以标准容器（如带时间戳的二进制或通用标记语言描述）存储，
并遵循监测机构的数据规范以便共享。
存储容量与回放工具决定“能追溯多久、
能否复盘某次干扰事件”，
是监管取证的关键。
理解监测站的体系形态与数据链，
才能正确使用监测接收机这一“监管工具”而非单纯“看谱仪器”。

\section{监测接收机基本原理}
\label{sec:monitoring_principle}

监测接收机的基本工作流程：

1. 
\textbf{频段扫描}：
按一定顺序扫描指定频段2. 
\textbf{信号检测}：
检测频段内的无线电信号3. 
\textbf{信号分析}：
分析信号的频率、
幅度、
调制方式等参数4. 
\textbf{数据记录}：
记录信号数据和分析结果5. 
\textbf{显示输出}：
显示监测结果

监测接收机通常采用超外差架构，
结合数字信号处理技术实现高性能信号分析。

监测接收机的传统实现是“超外差 + 步进扫描”：
本振按步进扫过关注频段，
每步将对应射频搬移到固定中频，
由检波器或 A/D 取得该频点的幅度，
连点成线即得到一幅频谱“全景图”（panorama）。
变频关系仍为

\[
f_{\mathrm{IF}}=\left| f_{\mathrm{LO}}-f_{\mathrm{RF}} 
\right|
\]

现代数字中频监测接收机则在宽带中频上对一段频谱一次采样，
用快速傅里叶变换（FFT）直接得到该段内的功率谱，
扫描变成“按块（segment）跳变”，
速度远高于逐点步进。
FFT 的频率分辨率为

\[
\Delta f=\frac{f_s}{N}
\]

其中 $f_s$ 为采样率，$N$ 为 FFT 点数；提高 $N$ 可细化频率分辨，但会增加一次变换所需时间，故分辨力与更新速率需折中。为兼顾“宽频段全景”与“细分辨力”，常用“全景扫描（粗、快）+ 感兴趣频段驻留分析（细、慢）”的两级策略。

信号分析环节对检出的信号做参数估计：
中心频率、
带宽、
电平（场强或功率）、
占用度，
以及调制识别（见下节）。
数据记录将频谱轨迹、
音频/中频样本与参数存入数据库，
供回放与统计（如频段占用率、
时间—频率—电平三维图）。
显示输出把上述结果以瀑布图、
频谱图或地图叠加方式呈现给操作员。

监测接收机的扫描策略有三种典型形态。
其一是传统步进扫描：
本振按频率步长逐点调谐，
每点驻留并检波，
连成频谱；
步长与驻留时间决定速度与精细度。
其二是扫频超外差：
本振连续扫过一段，
中频检波输出即随频率变化的包络，
适合宽频段快速“全景”，
但扫速过快会丢失瞬现信号。
其三是数字中频 FFT 块扫描：
对中频一段带宽一次采样做谱估计，
再跳到下一段，
兼顾速度与分辨，
是现代主流。

检波方式决定频谱轨迹的“形状”：
峰值检波保留最高电平、
适于发现瞬态；
平均值/RMS 检波反映平均功率、
适于占用度统计；
二者差别在宽带噪声下可达约 2.5dB（视带宽与视频带宽设置）。
视频带宽（VBW）用于平滑轨迹，
中频带宽（RBW）决定频率分辨——RBW 越窄分辨越细但扫描越慢。
瀑布图（时间—频率—电平的彩色图）则把连续频谱轨迹按时间堆叠，
便于发现间歇性、
跳频与突发信号。
信号参数估计在检测到信号后做：
定中心频率、
带宽、
电平，
并触发调制识别。

\section{监测接收机技术指标}
\label{sec:monitoring_specifications}

监测接收机的主要技术指标：

\begin{itemize}
    \item \textbf{频率范围}：能够覆盖的频率区间
    \item \textbf{灵敏度}：最小可检测信号强度
    \item \textbf{频率分辨率}：能够分辨的最小频率间隔
    \item \textbf{扫描速度}：单位时间内扫描的频段宽度
    \item \textbf{动态范围}：最大与最小可处理信号的比值
    \item \textbf{调制识别能力}：能够识别的调制方式类型
\end{itemize}

频率范围与灵敏度是基础指标：监测接收机常需覆盖数 kHz 至数 GHz 乃至更高，灵敏度（典型优于 $-120$ dBm 量级，配合预放可更低）决定能否发现微弱信号。频率分辨率指能区分两个相邻信号的最小间隔，数字中频中即由 $\Delta f=f_s/N$ 决定，硬件模拟检波时则受中频带宽限制。

扫描速度描述单位时间能“看”过多宽的频段。在步进扫描中，若每步驻留时间为 $\tau$、步长为 $\Delta f$，则平均扫描速率约 $\Delta f/\tau$；在数字中频按段采样时，完成一段带宽 $B$ 的一次谱估计所需时间约为 $T=N/f_s$，故实时带宽（可在不漏信号前提下连续分析的最大带宽）受 A/D 与处理吞吐限制。动态范围要求同时容纳强信号（如本地广播、雷达）与弱信号（如微小违规发射），通常用无杂散动态范围（SFDR）或无互调动态范围刻画，典型要求优于 70\sim80dB。

调制识别能力指自动判断信号属于 AM、
FM、
ASK、
FSK、
PSK、
QAM 等哪一类，
并估计符号率等参数，
是信号分类的关键。
下表汇总监测接收机典型指标的量级（数值为常见范围，
具体以设备指标书为准）。

\begin{table}[htbp]
\centering
\caption{监测接收机典型技术指标范围}
\begin{tabular}{|p{4.4cm}|p{5.6cm}|p{5.0cm}|}
\hline
\textbf{指标} & \textbf{含义} & \textbf{典型量级} \\
\hline
频率范围 & 可监测频段 & 数 kHz \sim 数 GHz \\
\hline
灵敏度 & 最小可检电平 & 优于 -120 dBm \\
\hline
频率分辨率 & 最小可分辨间隔 & 1Hz \sim 1kHz 可调 \\
\hline
扫描速度 & 单位时间扫过带宽 & 数十 MHz/s 量级 \\
\hline
动态范围 & 强弱信号并存能力 & 70\sim90 dB \\
\hline
调制识别 & 可识别调制类型数 & 十余种常见调制 \\
\hline
\end{tabular}
\end{table}

除前期列出的指标外，
监测接收机还有几项工程要点。
幅度精度与线性决定“测得的场强/功率是否可信”，
需依靠内部校准（功率检波器温补与自校准）与定期计量；
频率准确度由高稳参考（TCXO/OCXO 甚至 GPS 驯服）保证，
否则测得的信号频率会整体偏移。
相位噪声影响密近信号的分离与弱信号在强信号近旁的可见度；
镜频与中频抑制比决定带外强信号是否“虚假响应”到带内，
是宽频段监测的必检项。

截获概率（POI）与“实时带宽”是瞬现信号监测的关键：若扫描过快或处理跟不上，短时突发信号会被漏掉；现代设备以“在实时带宽内 100\% 截获最短脉宽”来表征。杂散响应（如 $m f_{\mathrm{LO}} \pm n f_{\mathrm{RF}}$ 落入中频）需通过频率规划与屏蔽抑制。综上，监测接收机指标是一组围绕“看得到、测得准、不漏、不乱报”的平衡：宽覆盖与高灵敏保证发现，高动态与低杂散保证可信，实时带宽与定位精度保证对瞬现与非法信号的处置能力。

频率参考决定“测得准不准”：监测接收机常用 TCXO 或 OCXO，关键固定站以 GPS 驯服钟获取长期稳定度，使频率读数偏差极小。幅度校准链决定“测得多真”：从射频端口到功率检波/数字功率，需经温补、自校准与定期计量，给出不确定度（如 $\pm1\sim2$ dB）。相位噪声影响密近信号的分离与弱信号在强信号近旁的可见度，镜频/中频抑制比决定带外强信号是否“虚假响应”到带内——二者是宽频段监测的必检项。

截获概率（POI）与实时带宽是瞬现信号监测的关键：扫描过快或处理跟不上会漏掉短时突发；现代设备以“在实时带宽内 100\% 截获最短脉宽”表征能力。杂散响应（如 $m f_{\mathrm{LO}} \pm n f_{\mathrm{RF}}$ 落入中频）需频率规划与屏蔽抑制，并在计量中逐项核查。综上，监测接收机指标围绕“看得到、测得准、不漏、不乱报”平衡：宽覆盖与高灵敏保发现，高动态与低杂散保可信，实时带宽与定位精度保对瞬现/非法信号的处置。

\section{监测接收机设计要点}
\label{sec:monitoring_design}

设计监测接收机时需考虑：

\begin{itemize}
    \item \textbf{宽带前端}：支持宽频段覆盖
    \item \textbf{高动态范围}：处理强弱信号
    \item \textbf{实时处理}：高速信号处理能力
    \item \textbf{存储能力}：大容量数据存储
    \item \textbf{远程控制}：支持远程监测和控制
\end{itemize}

\medskip\noindent\textbf{宽带前端与高动态范围}

监测接收机的射频前端要在极宽频段上保持平坦增益与低噪声，
同时承受可能出现的强信号而不致失真。
为此常采用“预选 + 分段放大 + 可调衰减”的结构：
在 broadband LNA 之前加可调预选器抑制带外强信号，
对极强调制信号自动插入衰减，
保护后级混频与 A/D。
混频器与 ADC 的三阶截获点（IP3）需足够高，
以保证在大信号背景下弱信号仍可被线性变换，
这是高动态范围的根本。

\medskip\noindent\textbf{实时处理、存储与远程控制}

数字中频要求 A/D 有足够采样率与有效位数（ENOB），
并有强大的 FPGA/DSP 完成连续 FFT、
数字滤波与检波，
以实现“实时带宽”内的无遗漏监测。
存储系统需保存长时间频谱轨迹与关键信号样本，
支持占用度统计与取证回放；
远程控制则通过标准总线（如网络、
SCPI 指令）把多台监测站组成分布式网络，
由中心站统一调度与汇总。
对移动监测车与便携设备，
还需兼顾功耗、
体积与抗振动。

\medskip\noindent\textbf{测向与信号分类识别}

测向（DF）定位是监测接收机的高级功能。单站测向常用干涉仪或到达角（AOA）法：以两天线构成基线 $d$，来波方向与基线夹角 $\theta$ 时两路接收信号的相位差为

\[
\Delta\varphi=\frac{2\pi d}{\lambda}\cos\theta
\]

由测得的多基线相位差可解算入射角；
多站则可用到达时间差（TDOA）做交叉定位，
精度取决于站间同步与时间测量。
信号分类识别则在参数估计基础上，
提取瞬时幅度、
频率、
相位以及谱特征（如峰均比、
谱对称性），
与模板库或分类模型比对，
自动给出调制体制、
带宽与可能用途的判断；
现代方法越来越多地引入机器学习提升对未知信号的识别率。

\section{监测接收机应用场景}
\label{sec:monitoring_applications}

监测接收机的主要应用：

\begin{itemize}
    \item \textbf{频谱管理}：监测频谱使用情况，合理分配频率资源
    \item \textbf{干扰监测}：检测和定位无线电干扰源
    \item \textbf{信号分析}：分析未知信号的特征和来源
    \item \textbf{合规监测}：检查无线电设备是否符合法规要求
    \item \textbf{安全监测}：监测非法无线电活动
\end{itemize}

在频谱管理方面，
监测接收机长期统计各频段占用率，
为频率规划、
指配与拍卖提供数据，
并发现“黑广播”“伪基站”等非法占用。
干扰监测结合测向定位，
可在民航、
铁路、
广电、
公众通信受扰时快速找到干扰源位置，
是电磁环境治理的关键手段。

信号分析面向未知或可疑信号，
提取其技术参数并研判属性，
用于技术侦察与态势感知；
合规监测则把被测台站的实际发射频率、
带宽、
功率、
占用度与执照参数比对，
判定是否越界或超标。
安全监测关注非法无线电活动（如违规遥控、
干扰、
未授权发射），
为无线电秩序与安全提供技术保障。
综上，
监测接收机是连接“频谱资源—台站—监管”的枢纽设备，
其数字化、
网络化与智能化正持续推进频谱管控能力的提升。

典型执法场景中，
“黑广播”（非法医药广告电台）与“伪基站”（冒充运营商发送诈骗短信）是常见查处对象：
监测站依频谱占用与调制特征发现异常发射，
结合测向与多站定位逼近源头，
由执法人员现场查封。
民航、
铁路专用频率受扰时，
监测车可沿线逼近定位干扰源（如违规大功率设备、
变频器泄漏），
保障运输安全。
重大活动与考点则部署临时监测，
防范 GPS 压制/欺骗等干扰。

频谱审计与规划是长期价值所在：
通过多年占用度数据识别闲置或低效频段，
支撑频率重整、
IMT 重耕与卫星轨位协调。
随着物联网与 5G/6G 设备激增，
电磁环境日益复杂，
监测接收机正向更宽实时带宽、
更强人工智能识别、
更低时延组网定位发展，
成为国家电磁空间安全与频谱资源高效利用不可或缺的技术基石。

监测工作须有法规与标准支撑：
国际电信联盟（ITU-R）对频谱监测方法、
场强与带宽定义给出建议，
各国无线电管理机构据此制定监测规范、
设备核准与干扰查处程序。
查处通常遵循“发现异常—参数测量—测向定位—现场核查—责令整改/查封”的流程，
监测数据是执法与技术判断的依据，
记录须完整、
可追溯以满足法律证据要求。

人员与体系建设同样关键：
监测操作员需理解电波传播、
接收机指标与测向误差来源，
避免把多径或站点误差误判为干扰位置；
跨区域与国际协调则在跨境干扰、
卫星干扰（如地球站受邻星干扰）中不可或缺。
面向未来，
人工智能将更多承担异常检测、
信号归并与趋势预测，
使有限的监测力量聚焦于真正重要的事件，
持续提升频谱治理的智能化与主动性。
